\documentclass[]{article}
\usepackage{amsmath}
\usepackage{amssymb}
\usepackage{mathtools}
\usepackage{graphicx}

\begin{document}
\title{Chernoff's Bound \& Hoeffding's Inequality}
\author{Soumyajit De}
\maketitle
\section{Why Care About Probability Bounds?}
Let's say that we have a random variable $X$ with an unknown distribution function $F$, i.e. $X\sim F$, which, typically, is the case. A common task is to estimate the true mean, $\mu=\mathbb{E}[X]$, assuming that it exists, from a few independent samples from the distribution, $X_1,\cdots,X_n\sim F$. We use the sample mean statistic, $\overline{X}_n=\frac{1}{n}\sum_{i=1}^n X_i$ as an estimator for this. This is reasonable, since the Weak-Law-of-Large-Numbers (WLLN) says that this quantity converges in-probability to the true-mean $\mu$ ($\overline{X}_n\overset{P}{\to} \mu$). Simply put, no matter how small a $t>0$ you choose, the probability of absolute difference between these two quantities being larger than that $t$ vanishing asymptotically, i.e. $\mathbb{P}(|\overline{X}_n-\mu|> t)\to 0$ as $n\to\infty$.
\\ \\ 
But in practice, when we have a finite sample $(X_i)_{i=1}^n$, this statement alone is not really useful if we want to quantify how far off this is from $\mu$. How quickly does it converge? How many samples would we need if we want the expected error in our estimate to be within a reasonable bound? 
\\ \\ 
Intuitively, these properties must depends on how concentrated the probability measure is around the mean. The more concentrated it is, the lesser number of samples we'd need to get a reasonable estimate. Ideally, we want the tails of the original distribution to be as thin as possible, which essentially means that the probability that the difference between $\mu$ and $X$ is large is bounded by a small quantity, i.e. $\mathbb{P}(|X-\mu|>t)\le\epsilon(t)$. Then, using these bounds, we can bound the error in our mean estimate $\overline{X}_n$ by some small quantity as well.
\\ \\
We'd see that the more information we have about the random variables and their distribution $F$, the tighter our bounds will be.
\section{Key idea: Use the information given}
\subsection{Only the mean is known}
\subsection{Both the mean and the variance are known}

\begin{center}
\includegraphics[scale=0.4]{notes/001.png}
\end{center}

\subsection{Variable is bounded}
\end{document}